\begin{scanpage}{0202}{184}
\[f(0)=f(\pi)=0,\tag{29.04}\]
which means that the radiation of a finite conductor in the direction of its geometric extension must be equal to zero. These conditions, together with a consideration of the relationships (28.05) and (28.16), lead to the system of equations.
\[
\left.\begin{aligned}
B_1+B_2\psi(L)e^{-2ikz_1}&=\psi(-z_1)e^{-2ikz_1},\\
B_1\psi(L)e^{2ikz_2}+B_2&=\psi(z_2)e^{2ikz_2}.
\end{aligned}\right\}\tag{29.05}
\]
from which we find without difficulty
\[
\left.\begin{aligned}
B_1&=\frac1D\bigl[\psi(-z_1)-\psi(L)\psi(z_2)e^{2ikz_2}\bigr]e^{-2ikz_1},\\
B_2&=\frac1D\bigl[\psi(z_2)-\psi(L)\psi(-z_1)e^{-2ikz_1}\bigr]e^{2ikz_2}.
\end{aligned}\right\}\tag{29.06}
\]
Keeping in mind (28.09), let us represent the functions $f(\vartheta)$ in a more graphic form
\[
\begin{aligned}
f(\vartheta)={}&1-\psi_+(z_2)e^{ikz_2(1-\cos\vartheta)}+\psi(z_2)e^{ik(L+z_2)}\times\\
&\times\bigl[\psi_-(L)e^{-ikz_1\cos\vartheta}-\psi(L)e^{ikL}\psi_+(L)e^{-ikz_2\cos\vartheta}\\
&\qquad+\psi^2(L)e^{2ikL}\psi_-(L)e^{-ikz_1\cos\vartheta}-\cdots\bigr]\\
&-\psi_-(-z_1)e^{-ikz_1(1+\cos\vartheta)}+\psi(-z_1)e^{ik(L-z_1)}\times\\
&\times\bigl[\psi_+(L)e^{-ikz_2\cos\vartheta}-\psi(L)e^{ikL}\psi_-(L)e^{-ikz_1\cos\vartheta}\\
&\qquad+\psi^2(L)e^{2ikL}\psi_+(L)e^{-ikz_2\cos\vartheta}-\cdots\bigr],
\end{aligned}\tag{29.07}
\]
where the secondary, tertiary, etc.\ waves corresponding to the first, second, and following terms in the brackets are explicitly written out.

Thus, the field radiated by a transmitting vibrator arises as a result of multiple diffraction of edge waves at the vibrator's ends. Let us note in connection with this that the edge wave is diffracted by the opposite end of the vibrator in the same way as at the end of
\end{scanpage}
